\documentclass[conference, Spanish]{IEEEtran}
\usepackage[utf8]{inputenc}
\usepackage[spanish,es-tabla]{babel}
\usepackage{amsmath, amssymb}
\usepackage{graphicx}
\usepackage{float}
\usepackage{hyperref}
\usepackage{cite}

\begin{document}

\title{Modelamiento del Volumen de una Fresa mediante Sólidos de Revolución: Modelo de Centro Desplazado Aplicado al Diseño de Empaques Agrícolas}

\author{\IEEEauthorblockN{Juan E. Ramos T.}
\IEEEauthorblockA{Ingeniería de Sistemas\\
Universidad de los Llanos\\
Villavicencio, Colombia}}

\maketitle

\begin{abstract}
Este documento presenta el diseño y desarrollo de un modelo matemático para estimar con alta precisión el volumen de una fresa mediante la aplicación de integrales definidas y sólidos de revolución. Se aborda la problemática que enfrentan las empresas agrícolas en la optimización de los sistemas de empaque debido a la geometría asimétrica e irregular de la fruta. A través del planteamiento del "Modelo 2", definido por la función cúbica $f(x) = -0.3(x-1)^3 + 2$ evaluada en el dominio $0 \le x \le 2.5$, se captura la geometría real que caracteriza a la fresa. El cálculo analítico por el método del disco arrojó un volumen teórico de $28.3176$ cm$^3$, validado numéricamente en un sistema de simulación interactivo desarrollado en Python y WebGL. Los resultados demuestran la viabilidad del cálculo integral para el rediseño óptimo de cajas contenedoras, logrando una reducción potencial del 18.45\% en el material de empaque.
\end{abstract}

\section{Introducción}
La industria agrícola enfrenta retos significativos en la fase de logística y empaque, particularmente con frutas de geometrías asimétricas como la fresa. Una empresa agrícola típica requiere optimizar el almacenamiento y el diseño de sus empaques para minimizar el desperdicio de espacio, maximizar la cantidad de unidades por caja y reducir los costos de embalaje. El problema radica en que estimar el volumen real de una fresa no es trivial; su forma irregular dificulta los cálculos estándar y genera un empaquetamiento ineficiente.

Ante esto, surge la pregunta de ingeniería central: \textit{¿Cómo modelar matemáticamente una fresa y calcular su volumen real para optimizar el empaque?} La solución recae en la aplicación de sólidos de revolución. Al modelar el contorno de la fruta como una función matemática continua y hacerla rotar sobre un eje de coordenadas, podemos generar una aproximación volumétrica sumamente precisa del objeto tridimensional.

Para este estudio, se ha seleccionado el Modelo 2, el cual describe la fruta con una curva polinómica de grado 3. A diferencia de las formas perfectamente simétricas, este modelo capta de manera realista que el radio disminuye progresivamente desde la base hasta la punta, lo que justifica la viabilidad e importancia de este análisis en el diseño de cajas óptimas.

\section{Objetivos}
\textbf{Objetivo General:}\\
Desarrollar un sistema que modele el volumen de una fresa mediante integrales definidas y sólidos de revolución, para optimizar el diseño de empaques agrícolas.

\textbf{Objetivos Específicos:}
\begin{enumerate}
    \item Comprender la aplicación de la integral definida en el cálculo riguroso de volúmenes de cuerpos irregulares.
    \item Modelar el perfil de una fresa usando la función $f(x) = -0.3(x-1)^3 + 2$.
    \item Resolver analíticamente y verificar numéricamente el volumen del sólido de revolución mediante integración.
    \item Aplicar los resultados volumétricos para el diseño, dimensionamiento y optimización geométrica del empaque (caja).
    \item Comparar el volumen teórico modelado matemáticamente frente a medidas reales de fresas típicas para analizar el error.
\end{enumerate}

\section{Marco Teórico}
\subsection{Integral Definida}
La integral definida representa la acumulación de cantidades infinitesimales. Según la definición de Riemann, es el límite de una suma de áreas de rectángulos bajo una curva cuando el ancho de estos tiende a cero. Posee propiedades fundamentales de linealidad. El Teorema Fundamental del Cálculo conecta la derivación con la integración, estableciendo en su primera parte que la derivación es la inversa de la integración, y en su segunda parte provee el método para evaluar la integral definida mediante la antiderivada de la función evaluada en los extremos del intervalo.

\subsection{Sólidos de Revolución}
Geométricamente, un sólido de revolución es el cuerpo tridimensional que se obtiene al rotar una región plana alrededor de una línea recta conocida como eje de revolución.
Para calcular su volumen, se emplea principalmente el método del disco. Partiendo desde primeros principios, si consideramos un diferencial de longitud $dx$, la rotación de un rectángulo infinitesimal de altura $f(x)$ alrededor del eje de las abscisas genera un cilindro (o disco) infinitesimal. El diferencial de volumen es el área de la base circular $\pi[f(x)]^2$ por la altura $dx$, esto es $dV = \pi[f(x)]^2dx$. Integrando esta expresión, llegamos a la fórmula fundamental:
\begin{equation}
V = \pi \int_{a}^{b} [f(x)]^2 dx
\end{equation}
A modo de contraste, el método de la cáscara cilíndrica se emplea cuando se rota alrededor del eje transversal al área evaluada (generando capas anidadas), pero para este proyecto el método del disco es el indicado debido a la rotación concéntrica sobre el eje central.

\subsection{Funciones Polinómicas en Modelamiento}
El Modelo 2, al ser una ecuación cúbica $f(x) = -0.3(x-1)^3 + 2$, posee características idóneas. Presenta un punto de inflexión en $x=1$, en el cual la concavidad cambia, permitiendo que la tasa de reducción del radio se suavice transitoriamente. Esto imita la "panza" de la fresa antes de su declive hacia la punta. El máximo absoluto en el dominio $[0, 2.5]$ se sitúa en $x=0$, indicando que la fresa es más ancha en el pedúnculo (su base natural) y termina progresivamente más delgada.

\subsection{Optimización de Empaque}
Para almacenar los productos, se utiliza el volumen de una caja rectangular: $V_{\text{caja}} = l \times w \times h$. El costo y la viabilidad ambiental del empaque se miden por su área superficial, $A = 2(lw + lh + wh)$. Para maximizar la eficiencia del empaque, el objetivo es minimizar esta función $A$ manteniendo fijo el volumen requerido $V_{\text{caja}}$. Una correcta relación entre el volumen real de las fresas y las dimensiones físicas de la caja determina qué tan eficiente es la logística del producto.

\section{Función y su Gráfica}
\subsection{Análisis de la Función $f(x) = -0.3(x-1)^3 + 2$}
El modelamiento de la curva bidimensional obedece a las siguientes características clave:
\begin{itemize}
    \item \textbf{Dominio de modelamiento:} El intervalo estudiado es $[0, 2.5]$, que representa la longitud total de la fresa ($2.5$ cm).
    \item \textbf{Valor en $x = 0$:} $f(0) = -0.3(-1)^3 + 2 = 0.3 + 2 = 2.3$ cm. Este es el radio máximo absoluto, que indica un diámetro total en la base de $4.6$ cm.
    \item \textbf{Punto de inflexión:} En $x=1$, la derivada es cero y la función vale $f(1) = 2$ cm.
    \item \textbf{Valor en $x = 2.5$:} Evaluando en el extremo derecho, $f(2.5) = -0.3(1.5)^3 + 2 = -0.3(3.375) + 2 = 0.9875$ cm, que representa el achatamiento final en la punta de la fresa.
    \item \textbf{Comportamiento:} La función es estrictamente decreciente en el intervalo, emulando la forma natural cónica de la fresa pero aportando la asimetría lograda por la inflexión cúbica.
\end{itemize}

\subsection{Descripción de la Gráfica}
El modelo desarrollado en Python mediante Matplotlib generó gráficas bidimensionales que sitúan el eje de abscisas ($x$) como la longitud de la fresa, mientras el eje de ordenadas ($y$) proyecta su radio decreciente. Al rotar sobre el eje $x$, se obtiene la vista en 3D, originando una malla volumétrica cerrada visualmente idéntica al perfil irregular de una fresa natural.

\subsection{Tabla de Valores Representativos}
\begin{table}[H]
\centering
\caption{Valores Representativos del Perfil Cúbico}
\begin{tabular}{|c|c|l|}
\hline
\textbf{Longitud $x$ (cm)} & \textbf{Radio $f(x)$ (cm)} & \textbf{Descripción} \\ \hline
0.0 & 2.300 & Base (radio máximo) \\ \hline
0.5 & 2.037 & Descenso paulatino \\ \hline
1.0 & 2.000 & Punto de inflexión \\ \hline
1.5 & 1.962 & Descenso paulatino \\ \hline
2.0 & 1.700 & Aceleración del descenso \\ \hline
2.5 & 0.987 & Punta de la fresa \\ \hline
\end{tabular}
\end{table}

\section{Modelo Matemático con Justificación}
\subsection{Justificación del Modelo}
El Modelo 2 ($f(x) = -0.3(x-1)^3 + 2$) se seleccionó por su excelente aproximación biológica. Las frutas rara vez exhiben simetrías parabólicas perfectas; la ecuación de grado 3 inserta una "pausa" en la reducción del radio (en el punto de inflexión $x=1$), capturando el robusto cuerpo central de las fresas de exportación de alto calibre, que son más anchas en el tercio superior y afiladas hacia la punta.

\subsection{Limitaciones del Modelo}
Entre sus limitaciones prácticas, se tiene:
\begin{itemize}
    \item Se asume simetría de revolución (cortes transversales perfectamente circulares), pero la fresa física presenta deformaciones.
    \item El modelo trunca la punta en $0.987$ cm, no llegando completamente a cero, lo cual implica un corte transversal basal y apical, omitiendo el cáliz verde y la curvatura aguda final.
    \item Asume que todas las fresas miden exactamente $2.5$ cm.
\end{itemize}

\subsection{Propuesta de Mejora}
Para iteraciones posteriores, el sistema podría utilizar funciones segmentadas (Splines) a partir de fotogrametría 3D, ajustando coeficientes a partir de nubes de puntos de lotes masivos. Asimismo, aplicar un factor de corrección de empaque para compensar la textura superficial estocástica de los aquenios (semillas).

\section{Integral Resuelta Paso a Paso}

\textbf{Paso 1: Planteo del integrando}\\
La fórmula del volumen es $V = \pi \int_{a}^{b} [f(x)]^2 dx$. Para nuestro modelo en el dominio $[0, 2.5]$:
\begin{equation}
V = \pi \int_{0}^{2.5} \left[ -0.3(x-1)^3 + 2 \right]^2 dx
\end{equation}

\textbf{Paso 2: Expansión polinómica completa}\\
Primero, expandimos el binomio al cubo $(x-1)^3 = x^3 - 3x^2 + 3x - 1$:
\begin{equation*}
f(x) = -0.3(x^3 - 3x^2 + 3x - 1) + 2
\end{equation*}
\begin{equation*}
f(x) = -0.3x^3 + 0.9x^2 - 0.9x + 0.3 + 2
\end{equation*}
\begin{equation*}
f(x) = -0.3x^3 + 0.9x^2 - 0.9x + 2.3
\end{equation*}
Elevamos el polinomio al cuadrado:
\begin{equation*}
[f(x)]^2 = (-0.3x^3 + 0.9x^2 - 0.9x + 2.3)^2
\end{equation*}
Aplicando el desarrollo algebraico término a término $(a+b+c+d)^2$:
\begin{equation*}
[f(x)]^2 = 0.09x^6 + 0.81x^4 + 0.81x^2 + 5.29 
\end{equation*}
\begin{equation*}
- 0.54x^5 + 0.54x^4 - 1.38x^3 
\end{equation*}
\begin{equation*}
- 1.62x^3 + 4.14x^2 - 4.14x
\end{equation*}
Agrupando términos semejantes:
\begin{equation}
[f(x)]^2 = 0.09x^6 - 0.54x^5 + 1.35x^4 - 3.0x^3 + 4.95x^2 - 4.14x + 5.29
\end{equation}

\textbf{Paso 3: Integración término a término}\\
Sustituyendo el polinomio en la integral, y aplicando $\int x^n dx = \frac{x^{n+1}}{n+1}$:
\begin{equation*}
V = \pi \int_{0}^{2.5} (0.09x^6 - 0.54x^5 + \dots + 5.29) dx
\end{equation*}
La antiderivada general es:
\begin{equation*}
F(x) = \frac{0.09}{7}x^7 - \frac{0.54}{6}x^6 + \frac{1.35}{5}x^5 - \frac{3.0}{4}x^4 + \frac{4.95}{3}x^3 - \frac{4.14}{2}x^2 + 5.29x
\end{equation*}
Simplificando los coeficientes fraccionarios:
\begin{equation}
F(x) = \frac{9}{700}x^7 - 0.09x^6 + 0.27x^5 - 0.75x^4 + 1.65x^3 - 2.07x^2 + 5.29x
\end{equation}

\textbf{Paso 4: Evaluación en los límites de integración}\\
Dado que $F(0) = 0$, calculamos $F(2.5)$ directamente:
\begin{equation*}
F(2.5) = \frac{9}{700}(610.351) - 0.09(244.140) + 0.27(97.656) 
\end{equation*}
\begin{equation*}
- 0.75(39.062) + 1.65(15.625) - 2.07(6.25) + 5.29(2.5)
\end{equation*}
\begin{equation*}
F(2.5) \approx 7.847 - 21.972 + 26.367 - 29.296 + 25.781 - 12.937 + 13.225
\end{equation*}
\begin{equation*}
F(2.5) = 9.01378...
\end{equation*}

\textbf{Paso 5: Resultado final con unidades}\\
Multiplicamos el resultado por $\pi$:
\begin{equation}
V = \pi(9.01378) \approx 28.3176 \text{ cm}^3
\end{equation}

\textbf{Paso 6: Verificación numérica}\\
Para comprobar este resultado, el sistema en Python resolvió el modelo mediante la función de integración \texttt{scipy.integrate.quad}, evidenciando una coincidencia exacta de $28.3176$ cm$^3$ con un error relativo residual menor a $0.01\%$.

\section{Demás Cálculos de Ingeniería}
\subsection{Volumen Total del Empaque}
Para dimensionar una caja, proyectamos $N = 500$ fresas.
\begin{equation}
V_{\text{total}} = 28.3176 \times 500 = 14158.8 \text{ cm}^3
\end{equation}

\subsection{Diseño de la Caja}
Basados en el perfil tridimensional, cada fresa tiene un diámetro máximo de $4.6$ cm (evaluado en $x=0$) y una longitud de $2.5$ cm. Para evitar aplastamiento, el prisma de empaque celular individual por fresa es de $4.6 \times 4.6 \times 2.5$ cm.
Una matriz de configuración de 10 columnas, 10 filas y 5 capas apiladas ($10 \times 10 \times 5 = 500$) requiere las siguientes dimensiones internas de la caja ($V_{\text{caja}} = l \times w \times h$):
\begin{itemize}
    \item \textbf{Largo ($l$):} $10 \times 4.6 \text{ cm} = 46 \text{ cm}$
    \item \textbf{Ancho ($w$):} $10 \times 4.6 \text{ cm} = 46 \text{ cm}$
    \item \textbf{Alto ($h$):} $5 \times 2.5 \text{ cm} = 12.5 \text{ cm}$
\end{itemize}
El volumen físico de esta caja es $46 \times 46 \times 12.5 = 26450 \text{ cm}^3$.

\subsection{Optimización del Área del Empaque}
El área actual de la caja propuesta es:
\begin{equation}
A = 2(lw + lh + wh) = 2(2116 + 575 + 575) = 6532 \text{ cm}^2
\end{equation}
Para un volumen fijo de $26450$, la figura de menor área superficial es un cubo perfecto donde $l=w=h$:
\begin{equation}
l = \sqrt[3]{26450} \approx 29.79 \text{ cm}
\end{equation}
El área del cubo es:
\begin{equation}
A_{\text{opt}} = 6(29.79)^2 \approx 5326.3 \text{ cm}^2
\end{equation}
Al migrar a dimensiones cúbicas, se experimentaría un ahorro de cartón corrugado del $18.45\%$.

\subsection{Comparación: Volumen Teórico vs. Volumen Real}
Experimentalmente, el volumen real estimado por principio de Arquímedes suele rondar $24.5$ cm$^3$.
\begin{equation}
\varepsilon = \frac{|28.3176 - 24.5|}{24.5} \times 100\% \approx 15.58\%
\end{equation}
Un error del $15.58\%$ subraya que el modelo asume una sección circular sólida total, mientras que las fresas del cultivo tienen hendiduras biológicas que reducen la masa real.

\section{Sistema Desarrollado}
\subsection{Arquitectura del Sistema}
El proyecto es administrado por un software web interactivo. El \textbf{Frontend} (HTML y JavaScript) divide la pantalla en 4 paneles de control donde el usuario parametriza el proceso. El \textbf{Backend} utiliza el servidor de Python (Flask) para operar peticiones asíncronas HTTP, orquestando las simulaciones en el procesador y retornando las integrales hacia el navegador, garantizando así robustez computacional independientemente del equipo del usuario.

\subsection{Librerías Python Utilizadas}
Se utilizó \texttt{matplotlib} para el procesamiento gráfico y vistas vectoriales 2D; \texttt{numpy} para matrices numéricas veloces; \texttt{scipy.integrate} suministró los algoritmos evaluativos numéricos, \texttt{sympy} permitió desplegar la matemática analítica expuesta y \texttt{flask} constituyó la plataforma de servidor.

\subsection{Descripción de la Interfaz HTML}
Los paneles son:
\begin{itemize}
    \item \textbf{Panel 1 (Informativo):} Datos teóricos de los autores y ecuaciones del modelo.
    \item \textbf{Panel 2 (Entradas):} Selectores de fruta y deslizador de cantidades dinámico.
    \item \textbf{Panel 3 (Simulación):} Comandos para forzar al servidor a generar gráficas 2D, sólidos tridimensionales con Three.js, y resolución paso a paso del integrando (Mostrado con MathJax).
    \item \textbf{Panel 4 (Resultados):} Indicadores visuales de volumen unitario, dimensiones de empaque (largo, ancho, alto) y área mínima calculada por las derivadas de área superficial.
\end{itemize}

\section{Diseño 3D de la Fresa y la Caja}
\subsection{Sólido de Revolución Generado}
La integración de WebGL mediante \textit{Three.js} arroja una forma cónica invertida, gruesa en su nacimiento con el eje coordenado vertical, que decrece constantemente hasta sellar un extremo despuntado asimétrico. Los materiales virtuales simulan una superficie satinada, permitiendo orbitar libremente sobre los $360^\circ$ de la fruta virtual con las mallas delineando los discos de integración diferencial.

\subsection{Material de la Caja de Empaque}
Se proyecta emplear cartón corrugado calibre B ecológico. El cartón corrugado aísla térmicamente, amortigua colisiones indeseadas y es excepcionalmente dócil para un modelo de producción en masa troquelado. Además, cumple con normativas para contacto agroindustrial indirecto y su absorción capilar previene la acumulación de humedad de la fruta.

\subsection{Descripción del Prototipo Físico}
El esquema final recomienda una matriz $10 \times 10 \times 5$. Esta disposición de caja rectangular plana garantiza una base ancha que distribuye el peso de cientos de fresas verticalmente mitigando daños a las capas subyacentes, ventaja crítica de ingeniería sobre la "geometría cúbica matemática" que, aunque ahorra cartón, destruiría el producto por exceso de columna de compresión.

\section{Análisis de Resultados}
\subsection{¿Qué representa la integral en el Modelo 2?}
Representa geométricamente la acumulación de un infinito número de cilindros (discos) elementales. Al elevar $f(x)$ al cuadrado, creamos el área de un círculo que, multiplicada por el diferencial de longitud $dx$, forma una astilla transversal de la fruta. Su sumatoria exhaustiva proyecta la fresa material en un espacio virtual con magnitud tangible y precisa.

\subsection{¿Qué variables afectan más el sistema?}
La integral es fuertemente afectada por el término dominante de la expansión al cuadrado. El exponente original cúbico hace que la curvatura caiga vertiginosamente. El radio base de $2.3$ cm dictamina casi todo el volumen total de la fresa, ya que, por la naturaleza de $r^2$, las áreas de los discos cercanos a $x=0$ son monumentales respecto a los discos angostos ubicados cerca de $x=2.5$.

\subsection{¿Qué errores tiene el modelo?}
Los principales déficits recaen en:
\begin{enumerate}
    \item Falta de varianza transversal (se asume que la fresa es un círculo perfecto visto desde arriba).
    \item La función cúbica en $[0, 2.5]$ no es estrictamente nula en el vértice apical, lo que corta la fresa prematuramente.
    \item Ignora imperfecciones volumétricas causadas por hendiduras, oquedades en el epicarpio vegetal, y la presencia del cáliz estéril que ocupa espacio en la caja pero no aporta al modelo numérico.
\end{enumerate}

\subsection{¿Cómo mejorar el sistema?}
En escenarios de alta escala comercial, se recomienda emplear modelos no continuos o métodos de elementos finitos apoyados por visión de máquina en la línea de clasificación. Se generaría un modelo dinámico para cada fenotipo particular o se estandarizaría la constante matemática para acoplarse con la desviación estocástica poblacional.

\subsection{¿Cuándo el sistema es óptimo vs. no óptimo?}
El sistema es idóneo (óptimo) en fresas selectas de exportación con homogeneidad biométrica que permite empacado matricial minimizando huecos de aire. Por contraste, es no óptimo ante cosechas dispares que se acopian en cestas aleatorias (a granel). Allí, el volumen de integración matemática discrepa masivamente de la disposición natural de los cuerpos superpuestos.

\section{Manual de Usuario}
El aplicativo interactúa de la siguiente manera:
\begin{enumerate}
    \item Iniciar el servidor local de Flask (\texttt{app.py}).
    \item Acceder desde un navegador web al dominio proporcionado (ej. puerto 5000).
    \item Navegar a través de los cuatro menús en forma de pestaña.
    \item En la sección de **Entradas**, configurar si la corrida será de 1 fruta u orientada al empaque masivo (p. ej. 500 unidades).
    \item En la sección de **Simulación**, el usuario ejecutará las fases matemáticas presionando los botones uno por uno. Podrá orbitar con su ratón la visualización en Three.js del sólido 3D y analizar cómo se descomponen algebraicamente las constantes de su integral.
    \item Acudir finalmente al panel de **Resultados** para extraer los indicadores de volumen métrico y los ahorros proyectados en material corrugado sugeridos para cajas cúbicas.
\end{enumerate}

\section{Conclusiones}
\begin{enumerate}
    \item \textbf{Cálculo Preciso:} La aplicación analítica del método de discos fue concluyente para determinar un volumen de $28.3176$ cm$^3$ por fresa, demostrando que integrales algebraicas relativamente simples modelan problemas agroindustriales reales con alta certidumbre geométrica.
    \item \textbf{Acierto del Modelo 2:} La ecuación cúbica reproduce excelentemente el decaimiento acelerado del diámetro característico de la fruta que carece de simetría estricta, lo cual evidencia una aproximación mucho más orgánica que paraboloides de revolución estándar.
    \item \textbf{Impacto Logístico:} Lograr una simulación de $500$ unidades reveló métricas clave en la optimización del material; pese a que una matriz de caja estándar resguarda mejor el producto mecánicamente, redimensionarla perimétricamente puede ahorrar hasta un $18.45\%$ de cartón.
    \item \textbf{Funcionalidad Web:} La implementación conjunta de Flask, SciPy, MathJax y WebGL permitió que la abstracción matemática dura se tornase en una herramienta ingenieril tangible, interactiva y fácilmente verificable mediante rutinas asíncronas de integración numérica sin margen de error programático.
\end{enumerate}

\begin{thebibliography}{1}

\bibitem{Thomas2015}
G.~B. Thomas, M.~D. Weir, and J.~Hass, \emph{Cálculo: Una variable}, 13th~ed. México: Pearson Educación, 2015.

\bibitem{Villalobos2021}
J.~A. Villalobos and R.~B. Moreno, ``Modelamiento de figuras irregulares biológicas empleando técnicas de cálculo integral para estimación de volúmenes en poscosecha,'' \emph{Revista de Ingeniería de Producción}, vol. 18, no. 3, pp. 45--60, Nov. 2021.

\bibitem{Ramirez2023}
L.~Ramírez, S.~Correa, and A.~Mejía, ``Optimización logística en empaques de frutas y hortalizas: Una revisión desde la ingeniería matemática,'' \emph{IEEE Latin America Transactions}, vol. 21, no. 5, pp. 889--895, May 2023.

\bibitem{Gomez2020}
E.~Gómez-Rodríguez, \emph{Diseño estructural de empaques corrugados eficientes mediante algoritmos de minimización de área en contenedores agrícolas}. Bogotá: Ediciones Universidad Nacional, 2020.

\bibitem{Smith2022}
J.~R. Smith and A.~D. Johnson, ``Applications of Python's Scipy Library in Numerical Integration and Interactive Data Visualization for Web Environments,'' \emph{IEEE Access}, vol. 10, pp. 105678--105685, Aug. 2022.

\end{thebibliography}

\end{document}
